\documentclass{article}
\usepackage[UTF8]{ctex}
\usepackage{amsmath,amssymb,amsfonts}
\usepackage{enumitem}
\usepackage{geometry}
\geometry{a4paper, margin=1in}
% === 错题本排版优化（不参与题目标识） ===
\setlength{\parindent}{0pt}
\setlength{\parskip}{0.35em plus 0.10em minus 0.08em}
\setlength{\emergencystretch}{3em}
\setlist{itemsep=0.25em, topsep=0.40em, parsep=0pt}
\allowdisplaybreaks[2]
% === 错题本排版优化结束 ===
\begin{document}
\section*{首都师范大学2026年数学分析真题}
\section*{33．首都师范大学2026年数学分析试卷}

\subsection{题目1}
一．（24 分）求极限．

\begin{itemize}
  \item[（1）] （8 分） \(\lim _{x \rightarrow 0} \frac{\tan x-x}{\sin ^{3} x}\) ．
  \item[（2）] （8 分） \(\lim _{n \rightarrow \infty} \sum_{k=1}^{n} \frac{k \sin \left(\frac{k}{n}\right)}{k n+1}\)
  \item[（3）] （8 分） \(\lim _{n \rightarrow \infty}\left(\frac{\sqrt[n]{a}+\sqrt[n]{b}}{2}\right)^{n}, a, b>0\) ．
\end{itemize}

\subsection{题目2}
二．（10 分）证明：\(f(x)=x \sin x\) 在 \([0,+\infty)\) 上不一致连续．

\subsection{题目3}
三。（15 分）解答如下问题：

\begin{itemize}
  \item[（1）] 证明：对任意给定的正整数 \(n\) ，方程 \(x^{n+1}+x^{n}+\cdots+x=1\) 在 \((0,1)\) 上存在唯一的根，记为 \(x_{n}\) ．
  \item[（2）] 证明：极限 \(\lim _{n \rightarrow \infty} x_{n}=\frac{1}{2}\) ．
\end{itemize}

\subsection{题目4}
四．（12 分）设

\[
f(x, y)= \begin{cases}x-y+\frac{x y^{2}}{x^{2}+y^{2}}, & (x, y) \neq(0,0) \\ 0, & (x, y)=(0,0)\end{cases}
\]

证明：\(f(x, y)\) 在原点 \((0,0)\) 处连续，沿任何方向的方向导数存在，但不可微．

\begin{itemize}
\end{itemize}

\subsection{题目5}
原题边界“五．”：（15 分）设 \(f(x)\) 在 \([0,1]\) 上连续可微，满足
\[
f^{\prime}(x)+f(x) \sin x=\int_{0}^{x} f(t) \mathrm{d} t, x \in[0,1]
\]
  \item[且] \(\int_{0}^{1} f(x) \mathrm{d} x=0\) ，证明：\(f(x) \equiv 0, x \in[0,1]\) ．

\subsection{题目5}
六．（10 分）证明：当 \(\lambda<1\) 时， \(\lim _{R \rightarrow+\infty} R^{\lambda} \int_{0}^{\frac{\pi}{2}} e^{-R \sin \theta} \mathrm{~d} \theta=0\) ．

\subsection{题目6}
七．（12 分）设 \(f(x)\) 在 \([-\pi, \pi]\) 上连续，且导函数 \(f^{\prime}(x)\) 可积和绝对可积，若有

\[
f(x) \sim \frac{a_{0}}{2}+\sum_{n=1}^{\infty}\left(a_{n} \cos n x+b_{n} \sin n x\right), x \in[-\pi, \pi] .
\]

证明：

\[
f^{\prime}(x) \sim \frac{c}{2}+\sum_{n=1}^{\infty}\left[\left(n b_{n}+(-1)^{n} c\right) \cos n x-n a_{n} \sin n x\right], x \in[-\pi, \pi] .
\]

其中 \(c=\frac{f(\pi)-f(-\pi)}{\pi}\) ，且 \(c=\lim _{n \rightarrow \infty}\left[(-1)^{n-1} n b_{n}\right]\) ．

\subsection{题目7}
八．（12 分）设 \(C\) 是单位圆周 \(x^{2}+y^{2}=1\) ，方向为逆时针，求积分

\[
\oint_{C} \frac{(x-y) \mathrm{d} x+(x+4 y) \mathrm{d} y}{x^{2}+4 y^{2}} .
\]

\subsection{题目8}
九．（10 分）计算曲面积分

\[
\iint_{S} x^{2} \mathrm{~d} y \mathrm{~d} z+y^{2} \mathrm{~d} z \mathrm{~d} x+z^{2} \mathrm{~d} x \mathrm{~d} y
\]

其中 \(S\) 为球面 \((x-a)^{2}+(y-b)^{2}+(z-c)^{2}=R^{2}(R>0)\) 的外侧．

\subsection{题目9}
十．（10 分）设函数列 \(f_{n}(x) \geq 0(n=1,2, \cdots)\) 与函数 \(g(x)\) 在区间 \([a, b]\) 上有界且黎曼可积，对任意的\(c \in(a, b)\) ，当 \(n \rightarrow \infty\) 时，\(\left\{f_{n}(x)\right\}\) 在区间 \([c, b]\) 上一致收敛于 0 ，且

\[
\lim _{n \rightarrow \infty} \int_{a}^{b} f_{n}(x) \mathrm{d} x=1, \lim _{x \rightarrow a^{+}} g(x)=A
\]

证明： \(\lim _{n \rightarrow \infty} \int_{a}^{b} g(x) f_{n}(x) \mathrm{d} x=A\) ．

\subsection{题目10}
十一．（10 分）设积分 \(\int_{0}^{+\infty} f(x) \mathrm{d} x\) 收敛，且存在 \(A>0\) ，使得 \(f(x)\) 在 \([A,+\infty)\) 上非负，存在 \(0<p<1\)使得当 \(x \rightarrow+\infty\) 时，\(x^{p} f(x)\) 单调递减趋于零，证明： \(\lim _{x \rightarrow+\infty} x f(x)=0\) ．

\subsection{题目11}
十二．（10 分）设 \(\left.f_{n}(x)\right\}(n=1,2, \cdots)\) 为 \([a, b]\) 上的连续函数列，且对任意给定的 \(x \in[a, b],\left\{f_{n}(x)\right\}\)关于 \(n\) 单调递增．证明：如果函数列 \(\left\{f_{n}(x)\right\}\) 在 \([a, b]\) 上逐点收敛于连续函数 \(f(x)\) ，则 \(\left\{f_{n}(x)\right\}\) 在\([a, b]\) 上必一致收敛于 \(f(x)\) 。
\end{document}
